\documentclass[11pt]{article}
\usepackage{amsmath,amssymb,amsthm}
\usepackage[margin=1in]{geometry}
\usepackage[colorlinks,citecolor=blue,linkcolor=blue]{hyperref}

\title{Nesting random regular graphs one degree at a time,\\ up to degree \(o(\sqrt{n})\)}
\author{AgenticOp-io}
\date{}

\newtheorem{theorem}{Theorem}
\newtheorem{lemma}[theorem]{Lemma}
\newtheorem{corollary}[theorem]{Corollary}
\newtheorem{remark}[theorem]{Remark}

\begin{document}
\maketitle

\begin{abstract}
Gao proved that \(G(n,d)\) and \(G(n,d+1)\) can be coupled so that
\(G(n,d)\subset G(n,d+1)\) a.a.s.\ whenever \(d\to\infty\) and \(d=o(n^{1/3})\).
We extend the range to \(d=o(n^{1/2}/\log^3 n)\) by substituting Gao's sharper
conditional edge probabilities into the same coupling template, replacing every
\(O(d^3/n)\) error by \(O(d^2/n)\).
\end{abstract}

\section{Introduction}

The Gao--Isaev--McKay monotonicity conjecture asks for couplings
\(G(n,d_1)\subset G(n,d_2)\) for \(d_1\le d_2\) (with two exceptional pairs).
The critical case is \(d_2=d_1+1\). Gao~\cite{Gao21} proved this for
\(d=o(n^{1/3})\) and observed that \(o(n^{1/2})\) should follow from the sharper
subgraph probabilities of~\cite{Gao20}. We carry out that program.

\begin{theorem}[Main]
\label{thm:main}
Assume the conditional edge-probability estimate of~\cite[Theorem~6]{Gao20}.
Let \(n\) be even and \(d\to\infty\) with \(d=o(n^{1/2}/\log^3 n)\).
Then there is a coupling of \(G_d\sim G(n,d)\) and \(G_{d+1}\sim G(n,d+1)\) with
\[
\mathbb{P}(G_d\subset G_{d+1})=1-o(1).
\]
\end{theorem}

\section{Patched subgraph probabilities}

\begin{lemma}[Patched Corollary~15]
\label{lem:cor15}
Let \(H\) be as in Gao~\cite[Corollary~15]{Gao21} with parameter \(\alpha=2k/n\).
If \(d=o(n^{1/2})\), then
\[
\mathbb{P}(H\subseteq G(n,d))
=(1+O(d^2/n))\,\rho(n,d,\alpha),
\]
with \(\rho\) the main term of~\cite{Gao21}.
\end{lemma}

\begin{proof}[Proof sketch]
Factor \(\mathbb{P}(H\subseteq G)\) as a product of conditional edge probabilities along
an ordering of \(E(H)\). Apply~\cite[Theorem~6]{Gao20} at each step (relative error
\(O(d^2/n^2)\)). Accumulate over \(|E(H)|\le n\) steps to get \(O(d^2/n)\).
The principal product equals Gao's main term by~\cite[Remark~8]{Gao20} together
with the error-free factorial calculation in~\cite{Gao21}.
\end{proof}

\begin{lemma}[Patched moments]
\label{lem:moments}
Let \(Y\) be the number of perfect matchings in \(G(n,d)\). Then
\[
\frac{\mathrm{Var}\,Y}{(\mathbb{E} Y)^2}
=\frac1{6d^3}+O\Bigl(d^{-4}+\frac{d^2}{n}+\sqrt{\frac{d}{n}}\log^3 n\Bigr).
\]
\end{lemma}

\begin{proof}[Proof sketch]
Identical to~\cite[Theorem~10]{Gao21}, with every \(O(d^3/n)\) from
Corollary~15 replaced by \(O(d^2/n)\) via Lemma~\ref{lem:cor15}.
\end{proof}

\section{Coupling}

Define
\[
\eta(\alpha)=2\alpha+\frac1{d^3\alpha^2}+C\frac{d^2}{n\alpha^2}
+C\frac{\sqrt{d/n}\,\log^3 n}{\alpha^2}.
\]
Choose
\[
\alpha=\max\bigl\{d^{-1},\,(d^2/n)^{1/4}\log n,\,(d/n)^{1/6}\log n\bigr\}.
\]
Then \(\alpha=o(1)\) and \(\eta(\alpha)=o(1)\) whenever
\(d=o(n^{1/2}/\log^3 n)\).

Run the coupling of~\cite[Theorem~19]{Gao21} with this \(\eta\) in place of Gao's
original error function. Chebyshev input comes from Lemma~\ref{lem:moments}.
Failure probability \(O(\eta)=o(1)\), proving Theorem~\ref{thm:main}.

\section{Remarks}

Constant-factor degree gaps for \(d\ge C\log n\) already follow from the
Kim--Vu sandwich theorem as proved in full by Behague--Il'kovi\v{c}--Montgomery.
The remaining barrier for the full monotonicity conjecture is
\(d=\Omega(\sqrt{n})\`.

\begin{thebibliography}{9}
\bibitem{Gao20}
P.~Gao,
Triangles and subgraph probabilities in random regular graphs,
arXiv:2012.01492, 2020.

\bibitem{Gao21}
P.~Gao,
The number of perfect matchings, and the nesting properties, of random regular graphs,
arXiv:2104.11850, 2021;
Random Structures \& Algorithms, 2023.
\end{thebibliography}

\end{document}
